% generated from Paper 6/anccft6.tex
\chapter{Algorithmic non-commutative class field theory, VI: the shadow algebra --- a functorial Kirchberg realization of the critical group}\label{chap:VI}
\chaptermark{VI}
\begin{chapterabstract}
The central open problem of this series --- posed in [Ch.~\ref{chap:III}, Rem.~1.10] and sharpened
in [Ch.~\ref{chap:IV}, \S6] --- asked for a Kirchberg $C^*$-algebra, functorial in the Hecke data of
a finite $(q+1)$-regular graph $Y$, whose $K_0$-group realizes the Laplacian-channel
invariant $\Z\oplus\Jac(Y)$, with $|\Jac(Y)|=\frac1n\prod_f\#E_f(\F_p)$ a product of
Frobenius point counts. Papers II--III proved two No-Go theorems that appeared to
surround the problem: no non-negative matrix of any size is shift equivalent to the
Ihara companion $C$ (the trace obstruction $\Tr(C^2)=-n(q-1)<0$), and the natural
tail-groupoid algebra is arithmetically blind. This paper solves the single-level
problem, inside the exact gap those theorems leave open: shift equivalence governs
spectra, not cokernels. Define the \emph{shadow graph} $\shadow$ on $2n$ vertices,
with adjacency
\[
B(Y)\;=\;\begin{pmatrix}T_p& I_n\\ qI_n & 2I_n\end{pmatrix}:
\]
the original graph, one arc from each vertex to its shadow, $q$ return arcs, and a
double loop at each shadow. The double loops are unimodular pivots of determinant
$-1$, and a two-line block factorization gives
$U\,(I-B^t)\,V=\Delta_p\oplus(-I_n)$ with explicit integer unimodular $U,V$; hence
the graph algebra $\cO_{B(Y)}$ is a unital Kirchberg algebra with
\[
K_0\bigl(\cO_{B(Y)}\bigr)\cong\Z\oplus\Jac(Y),\qquad
K_1\bigl(\cO_{B(Y)}\bigr)\cong\Z,
\]
functorially in $Y$. Being purely infinite it carries no trace, so the obstruction of
[Ch.~\ref{chap:II}] has no formulation on it; and it does not contradict [Ch.~\ref{chap:II}], because $B(Y)$ shares
with the Ihara pencil only its Bowen--Franks data, not its spectrum --- the exact
reconciliation is the per-eigenvalue identity
$(1-2u)(1-u\lambda)-qu^2\big|_{u=1}=-(q+1-\lambda)$. What remains open, and what [Ch.~\ref{chap:II}]
proves is impossible, are stated with equal precision: tower-level functoriality is
open; a \emph{Hecke-dynamical} realization is impossible, which is the sense in which
the present construction is maximal.
\end{chapterabstract}

\section*{Status of the results}

Theorems~\\ref{VI:thm:ktheory} and \\ref{VI:thm:kirchberg} are proved in full;
the unimodular factorization at the heart of Theorem~\\ref{VI:thm:ktheory} is a two-line
block computation reproduced below and machine-verified, together with the Smith-form
match against [Ch.~\ref{chap:I}, Table~1], on $K_4$, $K_5$, $K_{3,3}$ and Petersen
(Table~\\ref{VI:tab:battery}). The classical inputs --- simplicity and pure infiniteness
criteria for Cuntz--Krieger/graph algebras and their $K$-theory --- are quoted with a
flagged citation. Section~\\ref{VI:sec:scope} states exactly what is solved, what is
open, and what is impossible; one error in the task specification that prompted this
paper is corrected there (Remark~\\ref{VI:rem:k1}).

\section{The shadow graph}\label{VI:sec:shadow}

Throughout, $Y$ is a finite connected $(q+1)$-regular graph on $n$ vertices, $q\ge2$,
with adjacency (Hecke) matrix $T_p$ and Hecke--Laplacian $\Delta_p=(q+1)I-T_p$, whose
cokernel is $\Z\oplus\Jac(Y)$ with $\ker\Delta_p=\Z$ ([Ch.~\ref{chap:I}, Thm.~3.5]).

\begin{definition}[Shadow graph]\label{VI:def:shadow}
The \emph{shadow graph} $\shadow$ is the directed graph on vertex set
$V_Y\sqcup V_Y^{\,\prime}$ (each $v$ paired with its shadow $v'$) with arcs:
\begin{enumerate}
\item[(i)] $v\to w$ for every edge $vw$ of $Y$, in both orientations (the $T_p$
block);
\item[(ii)] one arc $v\to v'$ for every $v$;
\item[(iii)] $q$ parallel arcs $v'\to v$ for every $v$;
\item[(iv)] two loops at every shadow $v'$.
\end{enumerate}
Its adjacency matrix is $B(Y)=\begin{psmallmatrix}T_p&I\\qI&2I\end{psmallmatrix}$,
and the \emph{shadow algebra} is the graph $C^*$-algebra $\cO_{B(Y)}$.
\end{definition}

The design principle is worth stating in one sentence, because it is precisely the
loophole left by the No-Go theorems of [Ch.~\ref{chap:II}]. The obstruction
$\Tr(C^2)=-n(q-1)<0$ forbids any non-negative matrix from carrying the
\emph{spectrum} of the Ihara companion --- the $-qI$ block cannot be dilated away.
But $K$-theory sees only the \emph{cokernel}, which is a unimodular-equivalence
invariant, and unimodular row operations, unlike similarities, may be implemented by
pivots of determinant $-1$: a shadow vertex with a double loop contributes the
$1\times1$ block $1-2=-1$, and Schur elimination against it \emph{adds} $q$ back to
the diagonal --- exactly the operation that is spectrally impossible and
cohomologically free.

\section{The \texorpdfstring{$K$}{K}-theory}\label{VI:sec:ktheory}

\begin{theorem}[Functorial realization]\label{VI:thm:ktheory}
With $U=\begin{psmallmatrix}I&-qI\\0&I\end{psmallmatrix}$ and
$V=\begin{psmallmatrix}I&0\\-I&I\end{psmallmatrix}$, both integer unimodular,
\begin{equation}\label{VI:eq:factor}
U\,\bigl(I_{2n}-B(Y)^t\bigr)\,V\;=\;\Delta_p\oplus(-I_n).
\end{equation}
Consequently
\[
K_0\bigl(\cO_{B(Y)}\bigr)=\coker_\Z\bigl(I-B(Y)^t\bigr)\cong\Z\oplus\Jac(Y),
\qquad
K_1\bigl(\cO_{B(Y)}\bigr)=\ker_\Z\bigl(I-B(Y)^t\bigr)\cong\Z ,
\]
and the assignment $Y\mapsto\cO_{B(Y)}$ is functorial under graph isomorphisms: any
automorphism $\pi$ of $Y$ satisfies $(\pi\oplus\pi)B(Y)(\pi\oplus\pi)^{-1}=B(Y)$,
inducing an action on $\cO_{B(Y)}$ compatible with the induced action on
$\Z\oplus\Jac(Y)$.
\end{theorem}

\begin{proof}
$B^t=\begin{psmallmatrix}T_p&qI\\I&2I\end{psmallmatrix}$, so
$I-B^t=\begin{psmallmatrix}I-T_p&-qI\\-I&-I\end{psmallmatrix}$. Two lines of block
multiplication:
\[
U(I-B^t)=\begin{pmatrix}I&-qI\\0&I\end{pmatrix}
\begin{pmatrix}I-T_p&-qI\\-I&-I\end{pmatrix}
=\begin{pmatrix}(q+1)I-T_p&0\\-I&-I\end{pmatrix},
\]
\[
U(I-B^t)V=\begin{pmatrix}\Delta_p&0\\-I&-I\end{pmatrix}
\begin{pmatrix}I&0\\-I&I\end{pmatrix}
=\begin{pmatrix}\Delta_p&0\\0&-I\end{pmatrix}.
\]
Cokernel and kernel are unimodular-equivalence invariants, and
$\coker(\Delta_p\oplus(-I))=\coker\Delta_p=\Z\oplus\Jac(Y)$,
$\ker=\ker\Delta_p=\Z$ by [Ch.~\ref{chap:I}, Thm.~3.5(i)]. The $K$-theory formula for graph
$C^*$-algebras is the classical
$K_0=\coker(I-B^t)$, $K_1=\ker(I-B^t)$ \cite{CK}. Equivariance is the visible block
symmetry, verified explicitly for a transposition on $K_4$.
\end{proof}

\section{The Kirchberg properties}\label{VI:sec:kirchberg}

\begin{theorem}[Shadow algebra is Kirchberg]\label{VI:thm:kirchberg}
For $Y$ connected and $q\ge1$, the matrix $B(Y)$ is irreducible (the shadow graph is
strongly connected), is not a permutation matrix, and has a vertex --- any shadow
$v'$ --- carrying two distinct loops. Consequently $\cO_{B(Y)}$ is separable,
nuclear, simple, purely infinite and in the UCT class: a unital Kirchberg algebra.
Being purely infinite, it admits no nonzero trace, finite or semifinite; the trace
obstruction of {\rm[Ch.~\ref{chap:II}, Thms.~1.5--1.6]} therefore has no formulation on
$\cO_{B(Y)}$.
\end{theorem}

\begin{proof}
Strong connectivity: within the first layer all vertices are mutually reachable
through the arcs of the connected graph $Y$; layers communicate through
$v\to v'\to v$. The double loop at $v'$ gives two distinct return paths at one
vertex and period $1$. Irreducibility with a non-permutation matrix and the
double-loop witness place $B(Y)$ under the classical simplicity and pure
infiniteness criteria for Cuntz--Krieger algebras, and graph algebras of finite
graphs are separable, nuclear and satisfy the UCT \cite{CK}. Pure infiniteness
excludes traces. Machine confirmation of the witnesses on all four test graphs:
Table~\\ref{VI:tab:battery}.
\end{proof}

\section{Non-contradiction, and the exact scope of the solution}\label{VI:sec:scope}

\begin{proposition}[Why this does not contradict the No-Go theorems]\label{VI:prop:noncontra}
$B(Y)$ is not shift equivalent to the Ihara companion $C$ of {\rm[II]}: already
$\Tr B(Y)=2n\neq0=\Tr C$, and Lemma~{\rm[II,~1.4]} makes traces shift-equivalence
invariants. The shadow algebra shares with the Ihara pencil its Bowen--Franks data
and nothing spectral: per eigenvalue $\lambda$ of $T_p$,
\[
\det\bigl(I-uB(Y)\bigr)=\prod_{\lambda}\Bigl[(1-2u)(1-u\lambda)-qu^{2}\Bigr],
\qquad
(1-2u)(1-u\lambda)-qu^{2}\Big|_{u=1}=-(q+1-\lambda),
\]
so the $u=1$ specialization of the shadow zeta reproduces the Laplacian spectrum up
to sign, while for $u\neq1$ the two polynomial families differ. The construction
lives exactly in the gap the No-Go theorems left open: shift equivalence constrains
spectra; $K$-theory is a cokernel invariant; and Remark~[II, after Thm.~1.6]
anticipated precisely this distinction when it separated existence from canonicity.
\end{proposition}

\begin{proof}
The block determinant is computed with the commuting-block formula (three of the
four blocks are scalars), and the $u=1$ evaluation is
$(1-2)(1-\lambda)-q=\lambda-1-q$. Traces: immediate.
\end{proof}

\begin{remark}[A correction to the task specification]\label{VI:rem:k1}
The specification that prompted this paper asked for $K_1\cong\Z^{\,r}$. That is the
\emph{boundary}-channel invariant ([Ch.~\ref{chap:I}, Thm.~3.7]), not the Laplacian-channel one:
the target fixed by [Ch.~\ref{chap:I}, Thm.~3.5], [Ch.~\ref{chap:III}, Thm.~1.x] and [Ch.~\ref{chap:IV}] is
$K_1=\ker\Delta_p=\Z$, and that is what Theorem~\\ref{VI:thm:ktheory} delivers.
Realizing $\Z^r$ in $K_1$ alongside $\Z\oplus\Jac(Y)$ in $K_0$ would require a
non-square-type presentation and a different target; it is not the series' problem.
\end{remark}

\begin{remark}[What is solved, what is open, what is impossible]\label{VI:rem:scope}
\emph{Solved:} the single-level canonicity problem of [Ch.~\ref{chap:III}, Rem.~1.10] --- a
Kirchberg algebra functorial in the Hecke datum $T_p$, traceless, with
$K_0\cong\Z\oplus\Jac(Y)$, $K_1\cong\Z$, the torsion order
$\frac1n\prod_f\#E_f(\F_p)$ carried on the nose. \emph{Open:} tower-level
functoriality --- whether the tail norms of [Ch.~\ref{chap:V}, Thm.~3.1] lift to
$*$-homomorphisms or correspondences
$\cO_{B(Y_{k+1})}\to\cO_{B(Y_k)}$ inducing the norm maps of [IV--V] on $K_0$; nothing
here constructs such lifts. \emph{Impossible:} a \emph{Hecke-dynamical} realization,
i.e.\ an algebra whose gauge dynamics is the Ihara flow itself --- that is exactly
what [Ch.~\ref{chap:II}, Thm.~1.6] forbids for every non-negative presentation of every size. The
shadow algebra is Hecke-functorial, not Hecke-dynamical, and by [Ch.~\ref{chap:II}] no algebra can
be both: the construction is maximal in a theorem-certified sense.
\end{remark}

\begin{remark}[Normalization]\label{VI:rem:norm}
The construction uses the minimal unimodular pivot: loop weight $2$ is the least $m$
with $\det(1-m)=-1$, and the off-diagonal pair $(1,q)$ is one of the factorizations
$\alpha\beta=q$; any such choice yields the same $K$-theory by the same
factorization. We fix $(1,q,2)$ as the normal form and make no uniqueness claim
beyond it.
\end{remark}

\section{The verification battery}\label{VI:sec:battery}

\begin{table}[ht]
\centering\small
\renewcommand{\arraystretch}{1.2}
\begin{tabular}{lccccc}
\toprule
$Y$ & $n$ & $q$ & (S) factorization \\eqref{VI:eq:factor} & $\Jac(Y)$ from
$\mathrm{SNF}(I-B^t)$ & vs [Ch.~\ref{chap:I}, Table 1]\\
\midrule
$K_4$      & 4  & 2 & exact & $(\Z/4)^2$                 & match\\
$K_5$      & 5  & 3 & exact & $(\Z/5)^3$                 & match\\
$K_{3,3}$  & 6  & 2 & exact & $\Z/3\oplus\Z/3\oplus\Z/9$ & match\\
Petersen   & 10 & 2 & exact & $\Z/2\oplus(\Z/10)^3$      & match\\
\bottomrule
\end{tabular}
\medskip
\caption{Machine verification (\texttt{hecke\_kirchberg.py}): the explicit
factorization \\eqref{VI:eq:factor} holds as an exact matrix identity; the invariant
factors of $I-B(Y)^t$ reproduce [Ch.~\ref{chap:I}, Table~1] exactly; strong connectivity, the
double-loop witness, $\Tr B=2n\neq0$, and the $u{=}1$ reconciliation identity all
verified; automorphism equivariance spot-checked on $K_4$.}
\label{VI:tab:battery}
\end{table}

\section*{Acknowledgement of provenance and caveats}

Unrefereed preprint. Theorems~\\ref{VI:thm:ktheory}--\\ref{VI:thm:kirchberg} and
Proposition~\\ref{VI:prop:noncontra} are proved; the two-line factorization
\\eqref{VI:eq:factor} can be checked by hand and was machine-verified on four graphs.
The classical Cuntz--Krieger inputs (graph-algebra $K$-theory; simplicity and pure
infiniteness criteria) are quoted through \cite{CK}. Remark~\\ref{VI:rem:scope} delimits the result: single-level canonicity
solved, tower functoriality open, Hecke-dynamical realization impossible by [Ch.~\ref{chap:II}].
